%----------------------------------------------------------------------
%  Preamble
%----------------------------------------------------------------------
\usepackage[T1]{fontenc}
\usepackage[utf8]{inputenc}
% Body font. Swap freely: mathptmx (Times), charter, or lmodern if installed.
\usepackage{mathptmx}
\usepackage[scaled=0.92]{helvet}
\usepackage[english]{babel}
\usepackage{microtype}

% --- page layout ---
\usepackage[a4paper,left=3.0cm,right=2.5cm,top=2.5cm,bottom=2.5cm]{geometry}
\usepackage{setspace}
\onehalfspacing
\setlength{\parskip}{0.4em}

% --- maths ---
\usepackage{amsmath,amssymb,mathtools}
\usepackage{bm}

% --- tables and floats ---
\usepackage{booktabs}
\usepackage{tabularx}
\usepackage{longtable}
\usepackage{multirow}
\usepackage{threeparttable}
\usepackage{float}
\usepackage[font=small,labelfont=bf]{caption}
\usepackage{rotating}
\newcolumntype{L}{>{\raggedright\arraybackslash}X}

% --- lists, headers, colours ---
\usepackage{enumitem}
\setlist{itemsep=2pt,topsep=4pt}
\usepackage{xcolor}
\definecolor{linkblue}{RGB}{0,60,120}
\definecolor{phgray}{RGB}{120,120,120}
\usepackage{fancyhdr}
\pagestyle{fancy}
\fancyhf{}
\fancyhead[L]{\small\nouppercase{\leftmark}}
\fancyhead[R]{\small\thepage}
\renewcommand{\headrulewidth}{0.4pt}
\fancypagestyle{plain}{\fancyhf{}\fancyfoot[C]{\small\thepage}%
  \renewcommand{\headrulewidth}{0pt}}

% --- code listings ---
\usepackage{listings}
\lstset{
  basicstyle=\ttfamily\footnotesize,
  keywordstyle=\bfseries,
  commentstyle=\itshape\color{phgray},
  breaklines=true,
  frame=single,
  framerule=0.3pt,
  numbers=left,
  numberstyle=\tiny\color{phgray},
  captionpos=b,
  showstringspaces=false
}
% listings has no built-in Rust definition, so declare a minimal one.
\lstdefinelanguage{Rust}{
  keywords={pub, trait, fn, let, mut, struct, enum, impl, for, in, if,
            else, match, return, self, Self, use, mod, const, type,
            where, as, ref, Vec},
  keywordstyle=\bfseries,
  ndkeywords={u64, i64, f64, usize, bool, String, Option, Result},
  ndkeywordstyle=\itshape,
  sensitive=true,
  comment=[l]{//},
  morecomment=[s]{/*}{*/},
  morestring=[b]"
}
\lstset{language=Rust}

% --- citations ---
\usepackage[round,authoryear]{natbib}
\setcitestyle{aysep={},yysep={,}}

% --- hyperlinks (load late) ---
\usepackage[hidelinks,breaklinks=true]{hyperref}
\hypersetup{
  colorlinks=true,
  linkcolor=linkblue,
  citecolor=linkblue,
  urlcolor=linkblue,
  pdftitle={Frequent Batch Auction Engine: A Design Analysis and Simulation-Based Comparison against Continuous Trading},
  pdfsubject={Bachelor Thesis},
  pdfkeywords={market microstructure, frequent batch auction, matching engine, latency arbitrage, Hyperliquid}
}
\usepackage[capitalise,noabbrev]{cleveref}

% --- placeholder macros -------------------------------------------------
% \ph{} marks a value that the simulation has not yet produced. Every
% number in Chapter 5 that is not yet computed must use this macro, so
% that no placeholder can ever be mistaken for a measured result.
\newcommand{\ph}[1][\ensuremath{\cdot}]{\textcolor{phgray}{\textsf{[#1]}}}
\newcommand{\todo}[1]{\textcolor{red}{\textbf{[TODO: #1]}}}

% --- notation shorthands ---
\newcommand{\midp}{m_{t}}
\newcommand{\clr}{p^{\ast}}
\newcommand{\bps}{\,\text{bps}}
\DeclareMathOperator*{\argmax}{arg\,max}
\DeclareMathOperator*{\argmin}{arg\,min}

% --- theorem-like environments ---
\usepackage{amsthm}
\theoremstyle{definition}
\newtheorem{definition}{Definition}[chapter]
\newtheorem{proposition}{Proposition}[chapter]
\theoremstyle{remark}
\newtheorem{hypothesisx}{Prediction}[chapter]
\newenvironment{prediction}{\begin{hypothesisx}}{\end{hypothesisx}}

% --- chapter heading style ---
\usepackage{titlesec}
\titleformat{\chapter}[display]
  {\normalfont\huge\bfseries}{\chaptertitlename\ \thechapter}{16pt}{\Huge}
\titlespacing*{\chapter}{0pt}{-20pt}{28pt}

\widowpenalty=10000
\clubpenalty=10000
